\documentclass[openright,twoside,12pt]{report}

% Number subsubsections and include them in the TOC
\setcounter{secnumdepth}{3}
\setcounter{tocdepth}{3}

% PACKAGES
\usepackage{a4wide}
\usepackage{afterpage}
\usepackage{graphicx}
\usepackage[utf8]{inputenc}
\usepackage{mycxx}
\usepackage{myhyperref}
\usepackage{mylistings}
\usepackage{placeins}
\usepackage{supertabular}
\usepackage{xspace}

% SHORTHANDS
\newcommand{\ai}{AI\xspace}
\newcommand{\api}{API\xspace}
\newcommand{\caml}{Caml\xspace}
\newcommand{\io}{Io\xspace}
\newcommand{\java}{Java\xspace}
\newcommand{\js}{JavaScript\xspace}
\newcommand{\lua}{Lua}
\newcommand{\lisp}{Lisp}
\newcommand{\matlab}{MatLab\xspace}
\newcommand{\pascal}{PASCAL\xspace}
\newcommand{\python}{Python\xspace}
\newcommand{\ruby}{Ruby\xspace}
\newcommand{\sdk}{SDK\xspace}
\newcommand{\urbi}{Urbi\xspace}
\newcommand{\umake}{umake\xspace}
\newcommand{\us}{UrbiScript\xspace}

% FLOATS

% % This version tries to either put floats HERE, or regroup them on pages

% \newcommand   {\floatpos}          {htbp}
% \newcommand   {\floatposh}         {htbp}
% \newcommand   {\flushfloat}        {\afterpage{\clearpage}}
% \renewcommand {\floatpagefraction} {0.8}


% This version tries to insert float as near as possible from the
% declaration point, but after it.

% Don't insert a float at the top of a page before it's referred to!
\usepackage{flafter}

\newcommand   {\floatpos}          {htbp}
\newcommand   {\floatposh}         {!htb}
\newcommand   {\flushfloat}        {\afterpage{\clearpage}}
\setcounter   {topnumber}          {3} % Enable up to n floats at the top of pages
\setcounter   {bottomnumber}       {2} % Enable up to n floats at the bottom of pages
\renewcommand {\floatpagefraction} {0.8}


% LISTINGS

\lstset{language=[interactive]Urbi}

\definecolor{keywords}{rgb}{0.2,0.2,0.8}
\definecolor{comments}{rgb}{0.8,0.5,0}
\definecolor{string}{rgb}{0.0,0.5,0.0}
\definecolor{bg}{rgb}{0.9,0.9,1}

\lstset{
  keywordstyle=\color{keywords}\bfseries,
  commentstyle=\color{comments},
  stringstyle=\color{string},
  backgroundcolor=\color{bg},
  aboveskip=.6cm,
  captionpos=b,
  frame=single,
  basicstyle=\ttfamily\footnotesize,
  showstringspaces=false,
}

% COMMANDS
\newcommand{\lst}[1]{\autoref{lst:#1}}
\newcommand{\sect}[1]{\autoref{sect:#1}}
\newcommand{\tab}[1]{\autoref{tab:#1}}
\makeatletter
\renewcommand{\maketitle}{
  \begin{titlepage}

    \vfill

    \begin{figure}[htp]
      \centering
      \includegraphics[width=12cm]{img/urbi-logo}
    \end{figure}

    \begin{center}
      {\Huge\bf\@title\\}
      \vspace{1cm}
      {\Large \@author\\}
      \vspace{1cm}
      {\large \@date\\}
    \end{center}

    \vfill

    \begin{figure}[htp]
      \centering
      \includegraphics[width=12cm]{img/gostai}
    \end{figure}

    \vfill
    \global\let\title\relax
  \end{titlepage}
}
\makeatother

% STYLE
\hypersetup{colorlinks, citecolor=blue, linkcolor=blue, urlcolor=blue}

% CONFIGURATION
\title{\us 2.0 specifications}
\author{Kernel Team}

% DOCUMENT
\begin{document}

\maketitle
\tableofcontents

%FIXME: chapter is overkill
\chapter{Introduction}

\section{Status of this document}

This document defines the specifications of the \us language version
2.0. It defines the expected behavior from the \us interpreter, the
standard library, and the \sdk. It can be used to check that some code
is valid, or browse \us or \Cxx \api for a desired feature. Reading
some parts with no particular reason can also provide you with
advanced knowledge or subtleties about some \us aspects.

This document is however not fit to be used as a start guide to learn
\us, since it is not structured in a progressive manner and too
detailed. You can think of this document as a dictionary: one does not
% FIXME: give directions to the tutorials
learn a foreign language by reading a dictionary. The \us Tutorial, or
the live \us tutorial built in the interpreter are good introduction
to \us.

This document doesn't either aim at giving advanced programming
techniques. Its only goal is to defined how the language and its
libraries should behave.

\section{\us}

\us is an interpreted language designed for highly interactive and
parallel behavior programming. %FILLME

\section{Structuration of this document}

This document is divided in four chapters:

\begin{itemize}
\item Tools specifications: presentation and usage of the different
  tools available with the \urbi framework related to \us, such as the
  \urbi server, the command line client, \umake, \ldots
\item \us language specifications: definition of the core constructs
  of the language and their behavior.
\item \us standard library specifications: listing of all classes and
  methods provided in the standard library.
\item \urbi \sdk specifications: definition of the \urbi software
  development kit that enables you to interact with \urbi from \Cxx.
\end{itemize}

\chapter{Tools specifications}

\section{\urbi kernel}
\section{Console server}
\section{\umake}

\FloatBarrier
\chapter{\us language specifications}

\FloatBarrier
\section{Syntax}
\subsection{Comments}

Comments are used to document the code, and are totally ignored by the
\us interpreter. \Cxx-style and \C-style comments are supported.

\begin{itemize}
\item A \lstinline|//| introduces a comment that lasts until the end
  of the line.
\item A \lstinline|/*| introduces a comment that lasts until
  \lstinline|*/| is encountered. A difference with \C is that these
  comments nest; that is, if two \lstinline|/*| are encountered, the
  comment will end after two \lstinline|*/|, not one.
\end{itemize}

\lst{comments} illustrates this.

\begin{lstlisting}[caption=Comments,label=lst:comments,float=\floatpos]
  // C++ style comment
  /* C style comment */
  /* These comments /* do */ nest */
\end{lstlisting}

\subsection{Identifiers}
\label{sect:us-syn-id}

Identifiers in \us are composed of one or more alphanumeric or
underscore (\lstinline|_|) characters, given that they must not start
with a digit. That is, identifiers must match the
\lstinline|[a-zA-Z_][a-zA-Z0-9_]*| regular expression (see
\lst{identifiers}).  Additionally, identifiers must not match any of
the \us reserved word which are documented in \sect{syn-key}. The only
exception to this rule is \lstinline|new|, which can be used as the
method identifier in a method call, as shown in
\lst{identifiers-keywords}.  Finally, identifiers can also be written
between simple quotes (\lstinline|'|), in which case they may contain
any characters, as in \lst{escape-identifiers}.

\begin{lstlisting}[caption=Identifiers,label=lst:identifiers,float=\floatpos]
  var x;
  var foobar51;
  obj.met_hod();
  // Those are invalid.
// var 3x;
// obj.3x();
\end{lstlisting}

\begin{lstlisting}[caption=Identifiers cannot be keywords,
  label=lst:identifiers-keywords,float=\floatpos]
  // Those are invalid because "if" is a keyword.
// var if;
// obj.if();
  // However, keywords can be escaped with simple quotes
  var 'if';
  obj.'if'();
\end{lstlisting}

\begin{lstlisting}[caption=Escaping identifiers with simple quotes,
  label=lst:escape-identifiers,float=\floatpos]
  // Identifiers can be escaped with simple quotes
  var '#x';
  var '1 2 3';
  obj.'[]';
\end{lstlisting}

\subsection{Keywords}
\label{sect:syn-key}

Keywords are reserved words that cannot be used as identifiers, for
instance.  They are listed in \tab{keywords}.  Keywords can be
escaped, and thus used as identifier by surrounding them with simple
quotes (\lstinline|'|), as shown in \lst{escape-keyword}.

\renewcommand{\baselinestretch}{.85}
\begin{table}[\floatpos]
  \caption{Keywords}
  \label{tab:keywords}
  \centering
  \begin{tabular}{|c|c||c|c|}
    \hline
    Keyword                       & Remark                                  &
    Keyword                       & Remark                                  \\
    \hline
    \lstinline"and"               & Synonym for  operator                   &
    \lstinline"long"              & Unused                                  \\
    \lstinline"and_eq"            & Synonym for  operator                   &
    \lstinline"loop"              & \lstinline|loop&| and
                                    \lstinline-loop|- flavors               \\
    \lstinline"asm"               & Unused                                  &
    \lstinline"loopn"             & Deprecated in favor of for              \\
    \lstinline"at"                &                                         &
    \lstinline"mutable"           & Unused                                  \\
    \lstinline"auto"              & Unused                                  &
    \lstinline"namespace"         & Unused                                  \\
    \lstinline"bitand"            & Synonym for \lstinline|&| operator      &
    \lstinline"new"               &                                         \\
    \lstinline"bitor"             & Synonym for \lstinline-|- operator      &
    \lstinline"not"               & Synonym for \lstinline|!| operator      \\
    \lstinline"bool"              & Unused                                  &
    \lstinline"not_eq"            & Synonym for \lstinline|!=| operator     \\
    \lstinline"break"             &                                         &
    \lstinline"object"            &                                         \\
    \lstinline"call"              &                                         &
    \lstinline"onleave"           &                                         \\
    \lstinline"case"              &                                         &
    \lstinline"or"                & Synonym for \lstinline-||- operator     \\
    \lstinline"catch"             & Unused                                  &
    \lstinline"or_eq"             & Synonym for \lstinline-|=- operator     \\
    \lstinline"char"              & Unused                                  &
    \lstinline"private"           & Ignored                                 \\
    \lstinline"class"             &                                         &
    \lstinline"protected"         & Ignored                                 \\
    \lstinline"closure"           &                                         &
    \lstinline"public"            & Ignored                                 \\
    \lstinline"compl"             & Synonym for \lstinline|~|               &
    \lstinline"register"          & Unused                                  \\
    \lstinline"const"             & Unused                                  &
    \lstinline"reinterpret_cast"  & Unnused                                 \\
    \lstinline"const_cast"        & Unused                                  &
    \lstinline"return"            &                                         \\
    \lstinline"continue"          &                                         &
    \lstinline"short"             & Unused                                  \\
    \lstinline"default"           & Unused                                  &
    \lstinline"signed"            & Unused                                  \\
    \lstinline"delete"            &                                         &
    \lstinline"sizeof"            & Unused                                  \\
    \lstinline"do"                &                                         &
    \lstinline"static"            & Deprecated                              \\
    \lstinline"double"            & Unused                                  &
    \lstinline"static_cast"       & Unused                                  \\
    \lstinline"dynamic_cast"      & Unused                                  &
    \lstinline"stopif"            &                                         \\
    \lstinline"else"              &                                         &
    \lstinline"struct"            & Unused                                  \\
    \lstinline"emit"              &                                         &
    \lstinline"switch"            &                                         \\
    \lstinline"enum"              & Unused                                  &
    \lstinline"template"          & Unused                                  \\
    \lstinline"event"             &                                         &
    \lstinline"this"              &                                         \\
    \lstinline"every"             &                                         &
    \lstinline"throw"             & Unused                                  \\
    \lstinline"explicit"          & Unused                                  &
    \lstinline"timeout"           &                                         \\
    \lstinline"export"            & Unused                                  &
    \lstinline"try"               & Unused                                  \\
    \lstinline"extern"            & Unused                                  &
    \lstinline"typedef"           & Unused                                  \\
    \lstinline"external"          &                                         &
    \lstinline"typeid"            & Unused                                  \\
    \lstinline"float"             & Unused                                  &
    \lstinline"typename"          & Unused                                  \\
    \lstinline"for"               & \lstinline|for&| and                    &
                                    \lstinline-for|- flavors
    \lstinline"union"             & Unused                                  \\
    \lstinline"foreach"           & Deprecated in favor of for              &
    \lstinline"unsigned"          & Unused                                  \\
    \lstinline"freezeif"          &                                         &
    \lstinline"using"             & Unused                                  \\
    \lstinline"friend"            & Unused                                  &
    \lstinline"var"               &                                         \\
    \lstinline"from"              &                                         &
    \lstinline"virtual"           & Unused                                  \\
    \lstinline"function"          &                                         &
    \lstinline"volatile"          & Unused                                  \\
    \lstinline"goto"              & Unused                                  &
    \lstinline"waituntil"         &                                         \\
    \lstinline"if"                &                                         &
    \lstinline"wchar_t"           & Unused                                  \\
    \lstinline"in"                &                                         &
    \lstinline"whenever"          &                                         \\
    \lstinline"inline"            & Unused                                  &
    \lstinline"while"             & \lstinline|while&| and
                                    \lstinline-while|- flavors              \\
    \lstinline"int"               & Unused                                  &
    \lstinline"xor"               & Synonym for \lstinline|^| operator      \\
    \lstinline"internal"          & Deprecated                              &
    \lstinline"xor_eq"            & Synonym \lstinline|^=| operator         \\
    \hline
  \end{tabular}
\end{table}
\renewcommand{\baselinestretch}{1}

\begin{lstlisting}[caption=Escaping keywords,label=lst:escape-keyword,float=\floatpos]
  var 'if' = 21;
  [00000000] 21
  'if' * 2;
  [00000000] 42
\end{lstlisting}

\subsection{Literals}
\subsubsection{Durations}

Durations are floats (see \autoref{sect-syn-lit-flo}) followed by a
% FIXME: formulation
time unit. They are simply equivalent to the same float, expressed in
millisecond. For instance, \lstinline|1s|, which stands for ``one
second'', is strictly equivalent to \lstinline|1000|. Available units
and their equivalent are shown in \tab{durations}.

\begin{table}[\floatpos]
  \caption{Duration units}
  \label{tab:durations}
  \centering
  \begin{tabular}{|c|c|c|}
    \hline
    unit        & abbreviation & equivalence for n         \\
    \hline
    millisecond & ms           & $n$                       \\
    second      & s            & $n * 1000$                \\
    minute      & m            & $n * 1000 * 60$           \\
    hour        & h            & $n * 1000 * 60 * 60$      \\
    day         & d            & $n * 1000 * 60 * 60 * 24$ \\
    \hline
  \end{tabular}
\end{table}

\subsubsection{Floats}
\label{sect:us-syn-lit-float}

Literal floats consist in a succession of digit, representing the
integral part of the number in decimal base, possibly followed by a
dot (\lstinline|.|) and another succession of digit representing the
decimal part. Briefly, float literals must verify the
\lstinline|[0-9]+(\.[0-9]+)?| regular expression. \lst{literal-floats}
is an example of literal floats.

\begin{lstlisting}[caption=Literal floats,label=lst:literal-floats,float=\floatpos]
  1;
  [00000000] 1
  1.1;
  [00000000] 1.1
\end{lstlisting}

\subsubsection{Lists}
\label{sect:us-syn-lit-list}

Literal lists are represented with a comma-separated, potentially
empty list of arbitrary expressions enclosed in square brackets
(\lstinline|[]|), as shown in \lst{literal-lists}.

\begin{lstlisting}[caption=Literal lists,label=lst:literal-lists,float=\floatpos]
  []; // The empty list
  [00000000] []
  [1, 2, 3];
  [00000000] [1, 2, 3]
\end{lstlisting}

\subsubsection{Strings}

String are enclosed in double quotes (\lstinline|"|) and can contain
arbitrary characters, which stand for themselves, with the exception
of the escapement character, backslash (\lstinline|\|), as we can see
in \lst{literal-strings}. Backslash can introduce the following escapements:

\begin{itemize}
  \item \lstinline |\\|: an actual backslash
  \item \lstinline|\a|: a bell ring
  \item \lstinline|\b|: a backspace
  \item \lstinline|\f|: a form feed
  \item \lstinline|\n|: a line feed
  \item \lstinline|\r|: a carriage return
  \item \lstinline|\t|: a tabulation
  \item \lstinline|\v|: a vertical tabulation
\end{itemize}

Note that the interpreter print escaped strings. That is, line feed
will be printed out as \lstinline|\n| when a string result is
dumped. An actual line feed will of course be output if a string
content is printed with echo for instance.

\begin{lstlisting}[caption=Literal strings,label=lst:literal-strings,float=\floatpos]
  "";
  [00000000] ""
  "foo";
  [00000000] ""
  "a\nb"; // UrbiScript escapes string when dumping them
  [00000000] "a\nb"
  echo("a\nb"); // We can see there is an actual line feed
  [00000000] *** a
  ..b
  echo("a\\nb");
  [00000000] a\nb
\end{lstlisting}

\subsection{Operators}

\us supports many operators, most of which are inspired from
\Cxx. Their syntax is presented here, and they are sorted and
described with their original semantic - that is, \lstinline|+| is an
arithmetic operator that sums two numeric values. However, as in \Cxx,
these operators might be use for any other purpose - that is,
\lstinline|+| can also be used as the concatenation operator on lists
and strings. Their semantic is thus not limited to what is presented
here.

Tables in this section sort operators top-down, by precedence order.
Group of rows (not separated by horizontal lines) describe operators
that have the same precedence. Many operators are syntactic sugar that
bounce on a method. In this case, the equivalent desugared expression
is shown in the ``Equivalence'' column. This can help you determine
what method to override to define an operator for an object (see
\sect{fixme}).

Although code examples are present, this section does not intend do
define those operators behaviors; it only present their syntax and
grammatical rules, such as precedence and associativity. The behavior
of those operators is described in the documentation of the classes on
which they are applicable, in \autoref{chap:std}.

% Operator generators
\newcommand{\operatorhead}{Operator & Use & Associativity & Original semantic
  & Equivalence\\}


\newcommand{\operator}[6][ ]{\lstinline@#2@&\lstinline@#3@&#4&#5&\lstinline@#6@#1\\}
\newcommand{\boperator}[3]{\operator{#1}{a #1 b}{#2}{#3}{a.'#1'(b)}}
\newcommand{\poperator}[3]{\operator{#1}{#1a}{#2}{#3}{a.'#1'()}}

\newcommand{\operatordot}    {\operator  {.}    {a.b}              {-}     {Message sending}          {Not redefinable}       }
\newcommand{\operatordota}   {\operator  {.}    {a.b(args)}        {-}     {Message sending}          {Not redefinable}       }
\newcommand{\operatorsub}    {\operator  {[]}   {a[args]}          {-}     {Subscript}                {a.'[]'(args)}          }
\newcommand{\operatorsubass} {\operator  {[] =} {a[args] = v}      {-}     {Subscript assignment}     {a.'[]='(args, v)}      }
\newcommand{\operatorass}[2][ ]    {\operator[#1]
                                         {=}    {a = b}            {Right} {Assignment}               {updateSlot("a", b)}    }

\newcommand{\operatoriass}[1]{\operator  {#1=}  {a #1= b}          {Right} {In place assignment}      {a = a #1 b}            }
\newcommand{\operatorsiass}  {
    \operatoriass{+}
    \operatoriass{-}
    \operatoriass{*}
    \operatoriass{/}
    \operatoriass{\%}
    \operatoriass{\^}
    \operatoriass{\~}
}
\newcommand{\operatorinc}    {\operator  {++}   {a++}              {-}     {Incrementation}           {(a = a + 1) - 1}       }
\newcommand{\operatordec}    {\operator  {--}   {a--}              {-}     {Incrementation}           {(a = a - 1) + 1}       }

\newcommand{\operatoruplus}  {\poperator {+}    {-}                {Identity}               }
\newcommand{\operatorumin}   {\poperator {-}    {-}                {Opposite}               }
\newcommand{\operatorexp}    {\boperator {**}   {Right}            {Exponentiation}         }
\newcommand{\operatormult}   {\boperator {*}    {Left}             {Multiplication}         }
\newcommand{\operatordiv}    {\boperator {/}    {Left}             {Division}               }
\newcommand{\operatormod}    {\boperator {\%}   {Left}             {Modulo}                 }
\newcommand{\operatorplus}   {\boperator {+}    {Left}             {Sum}                    }
\newcommand{\operatorminus}  {\boperator {-}    {Left}             {Difference}             }
\newcommand{\operatorlshift} {\boperator {<<}   {Left}             {Left bit shift}         }
\newcommand{\operatorrshift} {\boperator {>>}   {Left}             {Right bit shift}        }
\newcommand{\operatoreq}     {\boperator {==}   {Non Associative}  {Equality}               }
\newcommand{\operatorneq}    {\boperator {!=}   {Non Associative}  {Inequality}             }
\newcommand{\operatorpeq}    {\boperator {===}  {Non Associative}  {Physical equality}      }
\newcommand{\operatorpneq}   {\boperator {!==}  {Non Associative}  {Physical Inequality}    }
\newcommand{\operatoraeq}    {\boperator {=~=}  {Non Associative}  {Approximative equality} }
\newcommand{\operatorinf}    {\boperator {<}    {Non Associative}  {Inferior}               }
\newcommand{\operatorinfeq}  {\boperator {<=}   {Non Associative}  {Inferior or equal}      }
\newcommand{\operatorsup}    {\boperator {>}    {Non Associative}  {Superior}               }
\newcommand{\operatorsupeq}  {\boperator {>=}   {Non Associative}  {Superior or equal}      }
\newcommand{\operatorbxor}   {\boperator {^}    {Left}             {Bitwise exclusive or}   }
\newcommand{\operatorneg}    {\poperator {!}    {Left}             {Logical negation}       }
\newcommand{\operatorand}    {\boperator {\&\&} {Left}             {Logical and}            }
\newcommand{\operatoror}     {\boperator {||}   {Left}             {Logical or}             }

\subsubsection{Arithmetic operators}

\us supports classic arithmetic operators, with the classic semantic
on numeric values. See \tab{arithmetic-operators} and
\lst{arithmetic-operators}.

\begin{table}[\floatposh]
  \caption{Arithmetic operators}
  \label{tab:arithmetic-operators}
  \centering
  \begin{tabular}{|c|c|c|c|c|c|}
    \hline
    \operatorhead
    \hline
    \operatoruplus
    \operatorumin
    \hline
    \operatorexp
    \hline
    \operatormult
    \operatordiv
    \operatormod
    \hline
    \operatorplus
    \operatorminus
    \hline
  \end{tabular}
\end{table}

\begin{lstlisting}[caption=Arithmetic operators,
  label=lst:arithmetic-operators,float=\floatposh]
  1 + 1;
  [00000000] 2
  1 - 2;
  [00000000] -1
  2 * 3;
  [00000000] 6
  10 / 2;
  [00000000] 5
  2 ** 10;
  [00000000] 1024
  -(1 + 2);
  [00000000] -3
\end{lstlisting}

\subsubsection{Assignment operators}

Assignment in \us can be performed with the \lstinline|=| operator. A
large variety of shorthands such as \lstinline|+=| exist. Those are
not directly redefinable since they are equivalent to a normal
affectation combined with another operator. See
\tab{assignment-operators} and \lst{assignment-operators}.


\begin{table}[\floatposh]
  \caption{Assignment operators}
  \label{tab:assignment-operators}
  \centering
  \begin{tabular}{|c|c|c|c|c|c|}
    \hline
    \operatorhead
    \hline
    \operatorass[\footnotemark]{}
    \operatorsiass
    \hline
  \end{tabular}
\end{table}
\footnotetext{For object fields only. Assignment to local variables
  cannot be redefined. }

\begin{lstlisting}[caption=Assignment operators,
  label=lst:assignment-operators,float=\floatposh]
  var x = 0;
  [00000000] 0
  x = 10;
  [00000000] 10
  x += 10;
  [00000000] 20
  x /= 5;
  [00000000] 4
  x %= 3;
  [00000000] 1
  x++;
  [00000000] 1
  x;
  [00000000] 2
\end{lstlisting}

\subsubsection{Bitwise operators}

\us supports some bitwise operators, although they are often use for
other purpose than bit-related operations. See \tab{bitwise-operators}
and \lst{bitwise-operators}.

\begin{table}[\floatposh]
  \caption{Bitwise operators}
  \label{tab:bitwise-operators}
  \centering
  \begin{tabular}{|c|c|c|c|c|c|}
    \hline
    \operatorhead
    \hline
    \operatorlshift
    \operatorrshift
    \hline
    \operatorbxor
    \hline
  \end{tabular}
\end{table}

\begin{lstlisting}[caption=Bitwise operators,
  label=lst:bitwise-operators,float=\floatposh]
  4 << 2;
  [00000000] 16
  4 >> 2;
  [00000000] 1
\end{lstlisting}

\subsubsection{Logic operators}

\us support classic logic operators. See \tab{logic-operators} and
\lst{logic-operators}.

\begin{table}[\floatposh]
  \caption{Logic operators}
  \label{tab:logic-operators}
  \centering
  \begin{tabular}{|c|c|c|c|c|c|}
    \hline
    \operatorhead
    \hline
    \operatorneg
    \hline
    \operatorand
    \hline
    \operatoror
    \hline
  \end{tabular}
\end{table}

\begin{lstlisting}[caption=Logic operators,
  label=lst:logic-operators,float=\floatposh]
  true && true;
  [00000000] true
  true || false;
  [00000000] true
  !true;
  [00000000] false
\end{lstlisting}

\subsubsection{Comparison operators}

\us support classic comparison operators. See \tab{comparison-operators} and
\lst{comparison-operators}.

\begin{table}[\floatposh]
  \caption{Comparison operators}
  \label{tab:comparison-operators}
  \centering
  \begin{tabular}{|c|c|c|c|c|c|}
    \hline
    \operatorhead
    \hline
    \operatoreq
    \operatorneq
    \operatorpeq
    \operatorpneq
    \operatoraeq
    \operatorinf
    \operatorinfeq
    \operatorsup
    \operatorsupeq
    \hline
  \end{tabular}
\end{table}

\begin{lstlisting}[caption=Comparison operators,
  label=lst:comparison-operators,float=\floatposh]
  0 < 0;
  [00000000] false
  0 <= 0;
  [00000000] true
  0 == 0;
  [00000000] true
  0 === 0;
  [00000000] false
  var x = 0;
  [00000000] 0
  x === x;
  [00000000] true
  x !== x;
  [00000000] false
\end{lstlisting}

\subsubsection{Miscellaneous operators}

Those operators cannot be classified in other categories. See
\tab{miscellaneous-operators} and \lst{miscellaneous-operators}. Note
that the subscript (square bracket) operator is variadic: it can take
any number of arguments, that will be passed to the \lstinline|'[]'|
method of the targeted object.

\begin{table}[\floatposh]
  \caption{Miscellaneous operators}
  \label{tab:miscellaneous-operators}
  \centering
  \begin{tabular}{|c|c|c|c|c|c|}
    \hline
    \operatorhead
    \hline
    \operatordot
    \operatordota
    \hline
    \operatorsub
    \operatorsubass
    \hline
  \end{tabular}
\end{table}

\begin{lstlisting}[caption=Miscellaneous operators,
  label=lst:miscellaneous-operators,float=\floatposh]
  0 < 0;
  [00000000] false
  0 <= 0;
  [00000000] true
  0 == 0;
  [00000000] true
  0 === 0;
  [00000000] false
  var x = 0;
  [00000000] 0
  x === x;
  [00000000] true
  x !== x;
  [00000000] false
\end{lstlisting}

\clearpage
\subsubsection{All operators summary}

\tab{operators-summary} is a summary of all operators, to highlight
the overall precedences. Operators are sorted vertically by decreasing
precedence. Groups of rows represent operators with the same
precedence.

\begin{table}[\floatposh]
  \caption{Operators summary}
  \label{tab:operators-summary}
  \centering
  \begin{tabular}{|c|c|c|c|c|c|}
    \hline
    Operator               & Use                                    & Associativity
    & Original semantic    & Equivalence                            \\
    \hline
    \operatordot
    \operatordota
    \hline
    \operatorsub
    \operatorsubass
    \hline
    \operatoruplus
    \operatorumin
    \hline
    \operatorexp
    \hline
    \operatormult
    \operatordiv
    \operatormod
    \hline
    \operatorplus
    \operatorminus
    \hline
    \operatorlshift
    \operatorrshift
    \hline
    \operatoreq
    \operatorneq
    \operatorpeq
    \operatorpneq
    \operatoraeq
    \operatorinf
    \operatorinfeq
    \operatorsup
    \operatorsupeq
    \hline
    \operatorbxor
    \hline
    \operatorneg
    \hline
    \operatorand
    \hline
    \operatoror
    \hline
    \operatorass
    \operatorsiass
    \hline
    \operatorinc
    \operatordec
    \hline
  \end{tabular}

\end{table}
\FloatBarrier

\FloatBarrier
\section{Scopes and local variables}

\subsection{Scopes}

% FIXME: This is wrong: the last separator isn't ignored; a pipe voids
% the result, for instance.
Scopes in \us are sequences of expressions, enclosed in curly brackets
(\lstinline|{}|). Expressions are separated with the four expression
separators (see \sect{fixme}). The last separator is optional and
ignored if present. Scopes are themselves expressions, and can thus be
used in composite expressions, nested, \ldots

\begin{lstlisting}[caption=Scopes,label=lst:scopes,float=\floatpos]
  // Scopes evaluate to their last expression
  {
    1;
    2;
    3;
  };
  [00000000] 3
  // Scopes can be used as expressions
  {1; 2; 3} + 1;
  [00000000] 4
\end{lstlisting}

\subsection{Local variables}

Local variables are introduced with the \lstinline|var| keyword,
followed by a valid identifier (see \sect{us-syn-id}) and an optional
initialization value assignment. If the initial value is omitted, it
defaults to void (\sect{std-void}). Variable declarations evaluate to
the initialization value. They can later be referred to by their
name. Their value can be changed with the assignment operator; such an
assignment expression returns the new value. The use of local
variables is illustrated in \lst{local-variables}.

\begin{lstlisting}[caption=Using local variables,
  label=lst:local-variables,float=\floatpos]
  // This declare variable x with value 42, and evaluates to 42.
  var x = 42;
  [00000000] 42
  // x equals 42
  x;
  [00000000] 42
  // We can assign it a new value
  x = 51;
  [00000000] 51
  x;
  [00000000] 51
  // Initialization defaults to void
  var y;
  y.isVoid;
  [00000000] true
\end{lstlisting}

Local variables lifespan is the same as their enclosing scope. They
are thus only accessible from their scope and its
subscopes\footnote{Local variables can actually escape their scope
  with lexical closures, as explained in \sect{us-fun-closures}}. Two
variables with the same name cannot be defined in the same scope. A
variable with the same name can be redefined in a subscope, in which
case references refer to the innermost variable, as shown in
\lst{local-variables-scoping}.

\begin{lstlisting}[caption=Local variables scoping,
  label=lst:local-variables-scoping,float=\floatpos]
  {
    var x = 0;
    var y = 1;
    {
      var y = 2;
      var z = 3;
      // We can access variables of parent scopes.
      echo(x);
      // This refers to the inner y.
      echo(y);
      echo(z);
    };
    // This refers to the outer y.
    echo(y);
    // This would be invalid: z does not exist anymore.
    // echo(z);
    // This would be invalid: x is already declared in this scope.
    // var x;
  };
  [00000000] *** 0
  [00000000] *** 2
  [00000000] *** 3
  [00000000] *** 1
\end{lstlisting}

\FloatBarrier
\section{Functions}

\subsection{Functions creation}

Functions in \us are first class citizen. Which mean, unlike \C and
related languages, a function is a normal value, just like an integer
or a string is. One can create a functional value thanks to the
\lstinline|function| keyword, followed by the list of formal arguments
and a scope representing the body of the function. Formal arguments
are a comma separated list of identifiers. \lst{function} illustrates
this.

\begin{lstlisting}[caption=Functional
  value,label=lst:function,float=\floatpos]
  function (arg1, arg2) { echo(0) };
  [00000000] function(arg1, arg2) {
  ..  echo(0)
  ..}
\end{lstlisting}

The whole point of a function is to be called later, so storing it in
a variable is the most common usage. \us provide a simple sugar to
initialize a new variable with a given function, presented in
\lst{function-sugar}.

\begin{lstlisting}[caption=Storing functions in variables,
  label=lst:function-sugar,float=\floatpos]
  // Functions are often stored in variables to call them later
  var f1 = function () {
    echo("hello")
  };
  f1();
  [00000000] *** hello
  // This form is strictly equivalent, yet simpler
  function f2()
  {
    echo("hello");
  }
  f2();
  [00000000] *** hello
\end{lstlisting}

\subsection{Arguments}

The list of formal arguments defines the number of argument the
function requires. They are accessible by their name from within the
body, as shown in \lst{function-args}. The list of formal arguments
can be omitted, in which case the number of given arguments isn't
check, and arguments themselves are not evaluated. Arguments can then
be manipulated with the call message, explained in
\sect{us-fun-callmsg}.

\begin{lstlisting}[caption=Storing and calling
  functions,label=lst:function-args,float=\floatpos]
  var f = function(a, b) {
    echo(b);
    echo(a);
  };
  f(1, 0);
  [00000000] *** 0
  [00000000] *** 1
  // Calling a function with the wrong number of argument is an error
  f(0);
  [00000000:error] !!! f: expected 2 arguments, given 1
  f(0, 1, 2);
  [00000000:error] !!! f: expected 2 arguments, given 3
\end{lstlisting}

\subsection{Return value}

The return value of the function is the evaluation of its body - that
is, since the body is a scope, the last evaluated expression in the
scope. However, value can be returned manually with the
\lstinline|return| keyword followed by the value, in which case the
evaluation of the function is stopped. If \lstinline|return| is used
with no value, the functions returns void. See \lst{function-return}.

\begin{lstlisting}[caption=Returning values from functions,
  label=lst:function-return,float=\floatpos]
  function f(a, b)
  {
    echo(a);
    echo(b);
    a // Return value is a
  };
  f(1, 2)
  [00000000] *** 1
  [00000000] *** 2
  [00000000] 1
  function g(a, b)
  {
    echo(a);
    return a; // Stop execution at this point and return a
    echo(b); // This is not executed
  };
  g(1, 2);
  [00000000] *** 1
  [00000000] 1
  function h()
  {
    return; // Stop execution at this point and return void
    echo(0); // This is not executed
  };
  h(); // Returns void, so nothing is printed.
\end{lstlisting}

\subsection{Call messages}
\label{sect:us-fun-callmsg}

Functions can access meta-information about how they were called,
through a \lstinline|CallMessage| object. The call message associated
with a function can be accessed with the \lstinline|call| keyword. It
contains several information such as not-yet evaluated arguments, the
name of the function, the target \ldots

One of the main usage of call message is to create lazy function, by
not specifying the formal arguments so as arguments are not evaluated
at the call site, and evaluate them by hand. A quick example is given
in \lst{function-lazy}. More information about this can be found in
the CallMessage class documentation (\sect{std-callmsg}).

\begin{lstlisting}[caption=Lazy
  function,label=lst:function-lazy,float=\floatposh]
  // Do not specify formal arguments, so as
  // effective arguments are not evaluated at call site
  function lazy
  {
    // Evaluate only the first argument, twice
    call.evalArgAt(0);
    call.evalArgAt(0);
  };
  lazy(echo("a"), echo("b"));
  [00000000] *** a
  [00000000] *** a
\end{lstlisting}

\subsection{Lexical closures}
\label{sect:us-fun-closures}

Lexical closure are an additional scoping rule, with which a function
can refer to a local variable located outside the function - but still
in the current context. \us supports read/write lexical closures,
meaning that the variable is shared between the function and the outer
environment, as shown in \lst{function-closure}.

\begin{lstlisting}[caption=Lexical
  closure,label=lst:function-closure,float=\floatpos]
  var x = 0;

  function cl()
  {
    // x refers to a variable outside the function
    x++;
    echo(x);
  };
  cl();
  [00000000] *** 1
  x;
  [00000000] 1
  x++;
  [00000000] 1
  cl();
  [00000000] *** 3
\end{lstlisting}

\lst{function-closure-escape} illustrate that local variables can even
escape their declaration scope by lexical closure.

\begin{lstlisting}[caption=Local variable escaping its scope by
  lexical closure.,label=lst:function-closure-escape,float=\floatpos]
  function wrapper()
  {
    // Normally, x is local to 'wrapper', and is limited to this scope.
    var x = 0;
    at (x > 1)
      echo("ping");
    // Here we make it escape the scope by returning a closure on it.
    return function() { x++ };
  };
  var f = wrapper();
  f();
  [00000000] 1
  f();
  [00000000] 2
  [00000000] *** ping
\end{lstlisting}

\FloatBarrier
\section{Objects}

Any value in \us is an object. Objects contain:

\begin{itemize}
\item A list of prototypes, which are also objects.
\item A list of slots, which to a name associate an object.
\end{itemize}

\subsection{Slots}

\subsubsection{Manipulation}

Objects can contain any number of slots, every slot has a name and a
value. They are often call ``fields'' or ``members'' in other object
oriented language.

The \lstinline|createSlot| is used to add a slot to an object with the
void (\sect{std-void}) value. The \lstinline|updateSlot| can be later
used to change the value of the slot, while the \lstinline|getSlot|
method can be used to read it. The \lstinline|setSlot| method can be
used to create a slot with a given value instead of void. Finally, the
\lstinline|slotNames| method can be used to get the list of the object
slot's name. \lst{slots} shows how to manipulate slots. More
documentation about these method can be found in \sect{std-object}.

\begin{lstlisting}[caption=Manipulating slots, label=lst:slots,
  float=\floatpos]
  var o = Object.new;
  o.slotNames;
  [00000000] []
  o.createSlot("test");
  [00000000] ["test"]
  o.slotNames;
  [00000000] ["test"]
  o.getSlot("test").asString;
  [00000000] "void"
  o.updateSlot("test", 42);
  [00000000] 42
  o.getSlot("test");
  [00000000] 42
\end{lstlisting}

\subsubsection{Sugars}

Several sugars exist to ease the use of slot methods:

\begin{itemize}
\item \lstinline|var o.name| is equivalent to
  \lstinline|o.createSlot("name")|.
\item \lstinline|var o.name = value| is equivalent to
  \lstinline|o.setSlot("name", value)|.
\item \lstinline|o.name = value| is equivalent to
  \lstinline|o.updateSlot("name", value)|.
\item \lstinline|delete o.name| is equivalent to
  \lstinline|o.removeSlot("name")|.
\end{itemize}


\subsection{Prototypes}

\subsubsection{Manipulation}

\us is a prototype-based language, unlike most classical object
oriented language, which are class-based. In prototype-based
languages, objects have no type, only prototypes, from which they
inherits behavior.

In \us, an object can have several prototypes. The list of prototypes
can be read with the \lstinline|protos| method. Prototypes can be add
or removed with the \lstinline|addProto| and \lstinline|removeProto|
methods respectively, as illustrated by \lst{prototypes}. More
documentation about these methods can be found in \sect{std-object}.

\begin{lstlisting}[caption=Manipulating prototypes, label=lst:prototypes,
  float=\floatpos]
  var o = Object.new;
  o.protos;
  [00000000] [<Object>]
  o.addProto(Pair);
  o.protos;
  [00000000] [<Object>, <Pair>]
  o.removeProto(Object);
  o.protos;
  [00000000] [<Pair>]
\end{lstlisting}

\subsubsection{Inheritance}

Objects inherit their prototypes' slots: \lstinline|getSlot| will also
look in an object prototypes' slots. \lstinline|getSlot| performs a
depth-first traversal of the prototypes hierarchy to find slots. That
is, when looking for a slot in an object:

\begin{itemize}
\item \lstinline|getSlot| checks first if the object itself has the
  requested slot. If so, it returns its value.
\item Otherwise, it applies the same research on every prototype, in
  the order of the prototype list (since addProto inserts in the front
  of the prototype list, the last prototype added has priority). This
  search is recursive: \lstinline|getSlot| will also look in the first
  prototype's prototype, etc before looking in the second
  prototype. If the slot is found in a prototype, it is returned.
\item Finally, if no prototype had the slot, an error is raised.
\end{itemize}

\lst{inheritance} shows how slots are inherited.

\begin{lstlisting}[caption=Slots inheritance, label=lst:inheritance,
  float=\floatpos]
  var a = Object.new;
  var b = Object.new;
  var c = Object.new;
  a.setSlot("x", "slot in a");
  b.setSlot("x", "slot in b");
  // c has no "x" slot
  c.getSlot("x");
  [00000000:error] !!! lookup failed: x
  // c can inherit the "x" slot from a
  c.addProto(a);
  c.getSlot("x");
  [00000000] "slot in a"
  // b is prepended to the prototype list, and has thus priority
  c.addProto(b);
  c.getSlot("x");
  [00000000] "slot in b"
  // a local slot in c has priority over prototypes
  c.setSlot("x", "slot in c");
  c.getSlot("x");
  [00000000] "slot in c"
\end{lstlisting}

\subsubsection{Copy on write}

The \lstinline|updateSlot| method has a particular behavior with
respect to prototypes. Although performing an \lstinline|updateSlot|
on a non existent slot is an error, it is valid if the slot is
inherited from a prototype. In this case, the slot is however not
updated in the prototype, but rather created in the object itself,
with the new value. This process is called {\em copy on write}; thanks
to it, prototypes are not altered when the slot is overridden in a
child object (see \lst{copy-on-write}).

\begin{lstlisting}[caption=Copy on write, label=lst:copy-on-write,
  float=\floatpos]
  var p = Object.new;
  var p.slot = 0;
  var c = Object.new;
  c.addProto(p);
  c.slot;
  [00000000] 0
  c.slot = 1;
  [00000000] 1
  // p's slot was not altered
  p.slot;
  [00000000] 0
  // It was copied in c
  c.slot;
  [00000000] 1
\end{lstlisting}

\subsection{Sending messages}

A message in \us consists in a message name and arguments. One can
send a message to an object with the dot (\lstinline|.|) operator,
followed by the message name (which can be any valid identifier) and
the arguments, as shown in \lst{messages}. When there are no
arguments, the empty parentheses can be omitted. As you might see,
sending messages is very similar to calling methods in classical
languages.

\begin{lstlisting}[caption=Sending messages, label=lst:messages,
  float=\floatpos]
  // Send the message msg to object obj, with arguments arg1 and arg2
  obj.msg(arg1, arg2);
  // Send the message msg to object obj, with no arguments
  obj.msg();
  // This is strictly equivalent
  obj.msg;
\end{lstlisting}

When a message ``msg'' is sent to object \lstinline|obj|:

\begin{itemize}
\item The ``msg'' slot of \lstinline|obj| is retrieved (i.e.,
  \lstinline|obj.getSlot("msg")|). If the slot is not found, the
  classic lookup error is raised.
\item If the object is not a \lstinline|Routine| (\sect{fixme}), it's
  the result of the message sending. In this case, there must be no
  argument, otherwise an error is raised.
\item If the object is a \lstinline|Routine|, it is invoked with the
  message sending arguments, and the returned value is the result. As
  a consequence, the number of arguments in the message sending must
  match the one required by the \lstinline|Routine|.
\end{itemize}

Such message sending is illustrated by \lst{send-message}.

\begin{lstlisting}[caption=Sending messages, label=lst:send-message,
  float=\floatpos]
  var obj = Object.new;
  var obj.a = 42;
  var obj.b = function (x) { x + 1 };
  obj.a;
  [00000000] 42
  obj.a();
  [00000000] 42
  obj.a(50);
  [00000000:error] !!! a: expected 0 arguments, given 1
  obj.b;
  [00000000:error] !!! b: expected 1 arguments, given 0
  obj.b();
  [00000000:error] !!! b: expected 1 arguments, given 0
  obj.b(50);
  [00000000] 51
\end{lstlisting}

\FloatBarrier
\section{Flow control}
\FloatBarrier
\section{Parallelism}
\FloatBarrier
\section{Events}
\FloatBarrier
\section{Pattern matching}
\FloatBarrier
\section{Trajectories}

\chapter{\us standard library specifications}
\label{chap:std}

\section{Boolean}
\section{CallMessage}
\label{sect:std-callmsg}
\subsection{Examples}
\subsubsection{Evaluating an argument several times}
\label{sect:std-callmsg-examples-several}
\section{Channel}
\section{Code}
\section{Comparable}
\label{sect:std-comparable}
\section{Dictionary}
\section{Event}

\section{Float}
\section{Float}

A Float is a floating point number.  It is also used, in the current
version of \us, to represent integers.

\subsection{Prototypes}

\begin{refObjects}
\item[Comparable]
\item[Orderable]
\item[RangeIterable]
\end{refObjects}

\subsection{Construction}
\label{sec:float:ctor}

The most common way to create fresh floats is using the literal
syntax.  Numbers are composed of three parts:
\begin{description}
\item[integral] (mandatory) a non empty sequence of (decimal) digits;
\item[fractional] (optional) a period, and a non empty sequence of
  (decimal) digits;
\item[exponent] (optional) either \samp{e} or \samp{E}, an optional
  sign (\samp{+} or \samp{-}), then a non-empty sequence of digits.
\end{description}

In other words, float literals match the
\lstinline|[0-9]+(\.[0-9]+)?([eE][-+]?[0-9]+)?|
regular expression.  For instance:

\begin{urbiassert}
0 == 0000.0000;
// This is actually a call to the unary '+'.
+1 == 1;
0.123456 == 123456 / 1000000;
1e3 == 1000;
1e-3 == 0.001;
1.234e3 == 1234;
\end{urbiassert}

There are also some special numbers, \lstinline|nan|, \lstinline|inf|
(see below).

\begin{urbiassert}
Math.log(0) == -inf;
Math.exp(-inf) == 0;
(inf/inf).asString == "nan";
\end{urbiassert}

A null float can also be obtained with \lstinline|Float|'s
\lstinline|new| method.

\begin{urbiassert}
Float.new == 0;
\end{urbiassert}

\subsection{Slots}

\begin{urbiscriptapi}
\item[abs]
  Absolute value of the target.
\begin{urbiassert}
(-5).abs == 5;
  0 .abs == 0;
  5 .abs == 5;
\end{urbiassert}

\item[acos]
  Arccosine of the target.
\begin{urbiassert}
0.acos == Float.pi/2;
1.acos == 0;
\end{urbiassert}

\item[asBool]
  Whether non null.
\begin{urbiassert}
0.asBool == false;
0.1.asBool == true;
(-0.1).asBool == true;
inf.asBool == true;
nan.asBool == true;
\end{urbiassert}

\item[asFloat]
  Return the target.
\begin{urbiassert}
51.asFloat == 51;
\end{urbiassert}

\item[asList]
  Bounces to \lstinline|seq|.  It is therefore possible to use the
  various flavors of \lstinline|for|-range loops on integers:
\begin{urbiassert}
{
  var res = [];
  for (var i : 3)
    res << i;
  res
}
== [0, 1, 2];

{
  var res = [];
  for|(var i : 3)
    res << i;
  res
}
== [0, 1, 2];

{
  var res = [];
  for&(var i : 3)
    res << i;
  res.sort
}
== [0, 1, 2];
\end{urbiassert}%>>

\item[asin]
  Arcsine of the target.
\begin{urbiassert}
0.asin == 0;
\end{urbiassert}

\item[asString]
  Return a string representing the target.
\begin{urbiassert}
42.asString == "42";
\end{urbiassert}

\item[atan]
  Return the arctangent of the target.
\begin{urbiassert}
0.atan == 0;
1.atan == Float.pi/4;
\end{urbiassert}

\item \lstinline|'bitand'(\var{that})|\\
  The bitwise-and between \lstinline|this| and \var{that}.
\begin{urbiassert}
(3 bitand 6) == 2;
\end{urbiassert}

\item \lstinline|'bitor'(\var{that})|\\
  Bitwise-or between \lstinline|this| and \var{that}.
\begin{urbiassert}
(3 bitor 6) == 7;
\end{urbiassert}

\item[clone]
  Return a fresh Float with the same value as the target.
\begin{urbiscript}
var x = 0;
[00000000] 0
var y = x.clone;
[00000000] 0
x === y;
[00000000] false
\end{urbiscript}

\item[compl]
  The complement to 1 of the target interpreted as a 32 bits integer.
\begin{urbiassert}
compl 0 == 4294967295;
compl 4294967295 == 0;
\end{urbiassert}

\item[cos]
  Cosine of the target.
\begin{urbiassert}
0.cos == 1;
Float.pi.cos == -1;
\end{urbiassert}

\item \lstinline|each(\var{fun})|\\
  Call the functional argument \var{fun} on every integer from 0 to
  target - 1, sequentially.  The number must be non-negative.
\begin{urbiassert}
{
  var res = [];
  3.each(function (i) { res << 100 + i });
  res
}
== [100, 101, 102];

{
  var res = [];
  for(var x : 3) { res << x; sleep(20ms); res << (100 + x); };
  res
}
== [0, 100, 1, 101, 2, 102];

{
  var res = [];
  0.each (function (i) { res << 100 + i });
  res
}
== [];
\end{urbiassert}

\item \lstinline'each|(\var{fun})'\\
  Call the functional argument \var{fun} on every integer from 0 to
  target - 1, with tight sequentiality.  The number must be
  non-negative.
\begin{urbiassert}
{
  var res = [];
  3.'each|'(function (i) { res << 100 + i });
  res
}
== [100, 101, 102];

{
  var res = [];
  for|(var x : 3) { res << x; sleep(20ms); res << (100 + x); };
  res
}
== [0, 100, 1, 101, 2, 102];
\end{urbiassert}%>>>>>>

\item \lstinline|each&(\var{fun})|\\
  Call the functional argument \var{fun} on every integer from 0 to
  target - 1, concurrently.  The number must be non-negative.
\begin{urbiassert}
{
  var res = [];
  for& (var x : 3) { res << x; sleep(30ms); res << (100 + x) };
  res
}
== [0, 1, 2, 100, 101, 102];
\end{urbiassert}%>>>>

\item[exp]
  Exponential of the target.
\begin{urbiscript}
1.exp;
[00000000] 2.71828
\end{urbiscript}

\item \lstinline|format(\var{finfo})|\\
  Format according to the \refObject{FormatInfo} object \var{finfo}.
  The precision, \lstinline|\var{finfo}.precision|, sets the maximum
  number of digits after decimal point when in fixed or scientific
  mode, and in total when in default mode.  Beware that 0 plays a
  special role, as it is not a ``significant'' digit.

  \begin{windows}
    Under Windows the behavior differs slightly.
  \end{windows}
\begin{urbiassert}
"%1.0d" % 0.1 == "0.1";
"%1.0d" % 1.1 == {if (System.Platform.isWindows) "1.1" else "1"};

"%1.0f" % 0.1 == "0";
"%1.0f" % 1.1 == "1";
\end{urbiassert}

\item[inf]
  Return the infinity.
\begin{urbiscript}
Float.inf;
[00000000] inf
\end{urbiscript}

\item[limit_digits]
  Number of digits (in \lstinline|Float.limit_radix| base) in the
  mantissa.
\begin{urbiassert}
Float.limit_digits;
\end{urbiassert}

\item[limit_digits10]
  Number of digits (in decimal base) that can be represented without
  change.
\begin{urbiassert}
Float.limit_digits10;
\end{urbiassert}

\item[limit_epsilon]
  Machine epsilon (the difference between 1 and the least value
  greater than 1 that is representable).
\begin{urbiassert}
1 != 1 + Float.limit_epsilon;
1 == 1 + Float.limit_epsilon / 2;
\end{urbiassert}

\item[limit_max]
  Maximum finite value.
\begin{urbiassert}
Float.limit_max     != Float.inf;
Float.limit_max * 2 == Float.inf;
\end{urbiassert}

\item[limit_max_exponent]
  Maximum integer value for the exponent that generates a normalized
  floating-point number.
\begin{urbiassert}
Float.inf != Float.limit_radix ** (Float.limit_max_exponent - 1);
Float.inf == Float.limit_radix ** Float.limit_max_exponent;
\end{urbiassert}

\item[limit_max_exponent10]
  Maximum integer value such that 10 raised to that power generates a
  normalized finite floating-point number.
\begin{urbiassert}
Float.inf != 10 ** Float.limit_max_exponent10;
Float.inf == 10 ** (Float.limit_max_exponent10 + 1);
\end{urbiassert}

\item[limit_min]
  Minimum positive normalized value.
\begin{urbiassert}
0 != Float.limit_min;
\end{urbiassert}

\item[limit_min_exponent]
  Minimum negative integer value for the exponent that generates a
  normalized floating-point number.
\begin{urbiassert}
0 != Float.limit_radix ** Float.limit_min_exponent;
\end{urbiassert}

\item[limit_min_exponent10]
  Minimum negative integer value such that 10 raised to that power
  generates a normalized floating-point number.
\begin{urbiassert}
0 != 10 ** Float.limit_min_exponent10;
\end{urbiassert}

\item[limit_radix]
  Base of the exponent of the representation.
\begin{urbiassert}
Float.limit_radix == 2;
\end{urbiassert}

\item[log]
  The logarithm of the target.
\begin{urbiassert}
0.log == -inf;
1.log == 0;
1.exp.log == 1;
\end{urbiassert}

\item \lstinline|max(\var{arg1}, ...)|\\
  Bounces to \lstinline|List.max| on \lstinline|[this, \var{arg1}, ...]|.
\begin{urbiassert}
1.max == 1;
1.max(2, 3) == 3;
3.max(1, 2) == 3;
\end{urbiassert}

\item \lstinline|min(\var{arg1}, ...)|\\
  Bounces to \lstinline|List.min| on \lstinline|[this, \var{arg1}, ...]|.
\begin{urbiassert}
1.min == 1;
1.min(2, 3) == 1;
3.min(1, 2) == 1;
\end{urbiassert}

\item[nan]
  The ``not a number'' special float value.  More precisely, this
  returns the ``quiet NaN'', i.e., it is propagated in the various
  computations, it does not raise exceptions.
\begin{urbiscript}
Float.nan;
[00000000] nan
(Float.nan + Float.nan) / (Float.nan - Float.nan);
[00000000] nan
\end{urbiscript}

A {NaN} has one distinctive property over the other Floats: it is
equal to no other float, not even itself.  This behavior is mandated
by the \wref[IEEE_754-2008]{IEEE 754-2008} standard.
\begin{urbiassert}
{ var n = Float.nan; n === n};
{ var n = Float.nan; n  != n};
\end{urbiassert}

\item[pi]
  $\pi$.
\begin{urbiassert}
Float.pi.cos ** 2 + Float.pi.sin ** 2 == 1;
\end{urbiassert}

\item[random]
  A random integer between 0 (included) and the target (excluded).
\begin{urbiscript}
20.map(function (dummy) { 5.random });
[00000000] [1, 2, 1, 3, 2, 3, 2, 2, 4, 4, 4, 1, 0, 0, 0, 3, 2, 4, 3, 2]
\end{urbiscript}

\item[round]
  The target, rounded to the nearest integer.
\begin{urbiassert}
1.6.round == 2;
1.4.round == 1;
\end{urbiassert}

\item[seq]
  The sequence of integers from 0 to \lstinline|this| - 1 as a list.
  The number must be non-negative.
\begin{urbiassert}
3.seq == [0, 1, 2];
0.seq == [];
(-1).seq;
[00004586:error] !!! seq: expected non-negative integer, got -1
\end{urbiassert}

\item[sign]
  Return 1 if \lstinline|this| is positive, 0 if it is null, -1
  otherwise.
\begin{urbiassert}
(-1164).sign == -1;
0.sign       == 0;
(1164).sign  == 1;
\end{urbiassert}

\item[sin]
  The sine of the target.
\begin{urbiassert}
0.sin == 0;
\end{urbiassert}

\item[sqr]
  Square of the target.
\begin{urbiassert}
32.sqr == 1024;
32.sqr == 32 ** 2;
\end{urbiassert}

\item[sqrt]
  The square root of the target.
\begin{urbiassert}
1024.sqrt == 32;
1024.sqrt == 1024 ** 0.5;
\end{urbiassert}

\item[srandom]
  Initialized the seed used by the random function.  As opposed to common
  usage, you should not use
\begin{urbiunchecked}
{
  var now = Date.now.timestamp;
  now.srandom;
  var list1 = 20.map(function (dummy) { 5.random });
  now.srandom;
  var list2 = 20.map(function (dummy) { 5.random });
  assert
  {
    list1 == list2;
  }
};
\end{urbiunchecked}

\item[tan]
  Tangent of the target.
\begin{urbiscript}
assert(0.tan == 0);
(Float.pi/4).tan;
[00000000] 1
\end{urbiscript}

\item \lstinline|times(\var{fun})|\\
  Call the functional argument \var{fun} \lstinline|this| times.

\begin{urbiscript}
3.times(function () { echo("ping")});
[00000000] *** ping
[00000000] *** ping
[00000000] *** ping
\end{urbiscript}

\item[trunc]
  Return the target truncated.
\begin{urbiassert}
1.9.trunc == 1;
(-1.9).trunc == -1;
\end{urbiassert}

\item \lstinline|'^'(\var{that})|\\
  Bitwise exclusive or between \lstinline|this| and \var{that}.
\begin{urbiassert}
(3 ^ 6) == 5;
\end{urbiassert}

\item \lstinline|'>>'(\var{that})|\\%>>
  \lstinline|this| shifted by \var{that} bits towards the right.
\begin{urbiassert}
4 >> 2 == 1;
\end{urbiassert}

\item \lstinline|'<'(\var{that})|\\
  Whether \lstinline|this| is less than \var{b}. The other comparison
  operators (\lstinline|<=|, \lstinline|>|, \ldots) can thus also be
  applied on floats since Float inherits \refObject{Orderable}.
\begin{urbiassert}
  0 < 1;
!(1 < 0);
\end{urbiassert}

\item \lstinline|'<<'(\var{that})|\\
  \lstinline|this| shifted by \var{that} bit towards the left.
\begin{urbiassert}
4 << 2 == 16;
\end{urbiassert}

\item \lstinline|'-'(\var{that})|\\
  \lstinline|this| subtracted by \var{b}.
\begin{urbiassert}
6 - 3 == 3;
\end{urbiassert}

\item \lstinline|'+'(\var{that})|\\
  The sum of \lstinline|this| and \var{that}.
\begin{urbiassert}
1 + 1 == 2;
\end{urbiassert}

\item \lstinline|'/'(\var{that})|\\
  The quotient of \lstinline|this| divided by \var{that}.
\begin{urbiassert}
50 / 10 == 5;
10 / 50 == 0.2;
\end{urbiassert}

\item \lstinline|'%'(\var{that})|\\
  \lstinline|this| modulo \var{b}.
\begin{urbiassert}
50 % 11 == 6;
\end{urbiassert}

\item \lstinline|'*'(\var{that})|\\
  Product of \lstinline|this| by \var{that}.
\begin{urbiassert}
2 * 3 == 6;
\end{urbiassert}

\item \lstinline|'**'(\var{that})|\\
  \lstinline|this| to the \var{that} power (${this}^{that}$).
\begin{urbiassert}
2 ** 10 == 1024;
\end{urbiassert}

\item \lstinline|'=='(\var{that})|\\
  Whether \lstinline|this| equals \var{that}.
\begin{urbiassert}
  1 == 1;
!(1 == 2);
\end{urbiassert}
\end{urbiscriptapi}

%%% Local Variables:
%%% mode: latex
%%% TeX-master: "../urbi-sdk"
%%% ispell-dictionary: "american"
%%% ispell-personal-dictionary: "../urbi.dict"
%%% End:


\section{Group}
\section{Integer}
\section{Lazy}

\subsection{Rationale}

Lazy are object that hold a lazy value; that is, a not-yet evaluated
value. They provide facilities to evaluate their content only once
(memoization) or several times. Lazy are essentially used in call
messages, to represent lazy arguments, as described in
\sect{std-callmsg}.

\subsection{Construction}

Lazy are seldom instantiated manually. They are mainly created
automatically when a lazy function call is made (see
\sect{us-fun-callmsg}). One can however create a lazy value with the
standard \lstinline|new| method of \lstinline|Lazy|, giving it an
argument-less function which evaluates to the lazified value
(\lst{new-lazy}).

\begin{lstlisting}[caption=Creating a lazy value, label=lst:new-lazy,
  float=\floatpos]
  Lazy.new(closure () { /* Value to lazify */ });
\end{lstlisting}

\subsection{Methods}
\subsubsection{eval}

The \lstinline|eval| method forces evaluation of the held lazy
value. Two calls to \lstinline|eval| will systematically evaluate the
expression twice, which can be useful to duplicates its side effects.

\subsubsection{value}

The \lstinline|value| method returns the held value, potentially
evaluating it before. \lstinline|value| performs memoization, that is,
only the first call will actually evaluate the expression, subsequent
calls will return the cached value. Unless you want to explicitly
trigger side effects from the expression by evaluating it several
time, this should be preferred over \lstinline|eval| to avoid
evaluating the expression several times uselessly.

\subsection{Examples}

\subsubsection{Evaluating once}

One usage of lazy values is to avoid evaluating an expression unless
it's actually needed, because it's expensive or has undesired side
effects. \lst{lazy-once} present a situation where an
expensive-to-compute value (\lstinline|heavy_computation|) might be
needed zero, one or two times. The objective is to save time by:

\begin{itemize}
\item Not evaluating it if it's not needed.
\item Evaluating it only once if it's needed one or two time.
\end{itemize}

We thus make the wanted expression lazy, and use the \lstinline|value|
method to fetch its value when needed.

\begin{lstlisting}[caption=, label=lst:lazy-once, float=\floatpos]
  // This function supposedly performs expensive computations.
  function heavy_computation()
  {
    echo("Heavy computation");
    return 1 + 1;
  };
  // We want to do the heavy computations only if needed,
  // and make it a lazy value to be able to evaluate it "on demand".
  var v = Lazy.new(closure () { heavy_computation() });
  [00000000] Lazy_0xADDR
  /* some code */;
  // So far, the value was not needed, and heavy_computation
  // was not evaluated.
  /* some code */;
  // If the value is needed, heavy_computation is evaluated.
  v.value;
  [00000000] *** Heavy computation
  [00000000] 2
  // If the value is needed a second time, heavy_computation
  // is not reevaluated.
  v.value;
  [00000000] 2
\end{lstlisting}

\subsubsection{Evaluating several time}

Evaluating a lazy several time only makes sense with lazy arguments
and call messages. See example with call messages in
\sect{std-callmsg-examples-several}.

\section{List}
\section{List}

\lstinline|List|s implement potentially-empty ordered (heterogeneous)
collections of elements.

\subsection{Prototypes}

\begin{itemize}
\item \refObject{Object}
\item \refObject{RangeIterable}
\item \refObject{Orderable}
\end{itemize}

\subsection{Construction}

List can be created with their literal syntax: a possibly empty
sequence of expressions in square brackets, separated by commas.
Non-empty list may actually \emph{terminate} with a comma, rather than
\emph{separate}; in other words, an optional trailing comma is accepted.

\begin{urbiscript}
[]; // The empty list
[00000000] []
[1, "2", [3,],];
[00000000] [1, "2", [3]]
\end{urbiscript}

\subsection{Slots}

\begin{itemize}
\item \lstinline|all(\var{fun})|\\
  % FIXME: link to predicate glossary entry
  Return whether all the members of the target verify the predicate
  \var{fun}.

\begin{urbiassert}[firstnumber=last]
// Are all elements positive?
! [-2, 0, 2, 4].all(function (e) { e > 0 });
// Are all elements even?
[-2, 0, 2, 4].all(function (e) { e % 2 == 0 });
\end{urbiassert}

\item \lstinline|any(\var{fun})|\\
  % FIXME: link to predicate glossary entry
  Whether at least one of the members of the target verifies the
  predicate \var{fun}.

\begin{urbiassert}[firstnumber=last]
// Is there any even element?
! [-3, 1, -1].any(function (e) { e % 2 == 0 });
// Is there any positive element?
[-3, 1, -1].any(function (e) { e > 0 });
\end{urbiassert}

\item \lstinline|asBool|\\
  Whether not empty.
\begin{urbiassert}[firstnumber=last]
[].asBool == false;
[1].asBool == true;
\end{urbiassert}

\item \lstinline|asList|\\
Return the target.

\begin{urbiassert}[firstnumber=last]
[0, 1, 2].asList == [0, 1, 2];
\end{urbiassert}

\item \lstinline|asString|\\
  A string describing the list.  Uses \lstinline|asPrintable| on its
  members, so that, for instance, strings are displayed with quotes.

\begin{urbiassert}[firstnumber=last]
[0, [1], "2"].asString == "[0, [1], \"2\"]";
\end{urbiassert}

\item \lstinline|back|\\
Return the last element of the target. An error if the target is empty.

\begin{urbiscript}[firstnumber=last]
assert([0, 1, 2].back == 2);
[].back;
[00000000:error] !!! back: cannot be applied onto empty list
\end{urbiscript}

\item \lstinline|clear|\\
  Empty the target.

\begin{urbiscript}[firstnumber=last]
var x = [0, 1, 2];
[00000000] [0, 1, 2]
assert(x.clear == []);
\end{urbiscript}

\item \lstinline|each(\var{fun})|\\
  Apply the given functional value \var{fun} on all members,
  sequentially.

\begin{urbiscript}[firstnumber=last]
[0, 1, 2].each(function (v) {echo (v * v); echo (v * v)});
[00000000] *** 0
[00000000] *** 0
[00000000] *** 1
[00000000] *** 1
[00000000] *** 4
[00000000] *** 4
\end{urbiscript}

\item \lstinline|'each&'(\var{fun})|\\
Apply the given functional value on all members simultaneously.

\begin{urbiscript}[firstnumber=last]
[0, 1, 2].'each&'(function (v) {echo (v * v); echo (v * v)});
[00000000] *** 0
[00000000] *** 1
[00000000] *** 4
[00000000] *** 0
[00000000] *** 1
[00000000] *** 4
\end{urbiscript}

\item \lstinline|empty|\\
  Whether the target is empty.

\begin{urbiassert}[firstnumber=last]
[].empty;
! [1].empty;
\end{urbiassert}

\item \lstinline|filter(\var{fun})|\\
  The list of all the members of the target that verify the predicate
  \var{fun}.

\begin{urbiassert}[firstnumber=last]
// Keep only odd numbers.
[0, 1, 2, 3, 4, 5].filter(function (v) {v % 2 == 1}) == [1, 3, 5];
\end{urbiassert}

\item \lstinline|foldl(\var{action}, \var{value})|\\
  \wref[Fold_(higher-order_function)]{Fold},
  also known as \dfn{reduce} or \dfn{accumulate}, computes a result
  from a list.  Starting from \var{value} as the initial result, apply
  repeatedly the binary \var{action} to the current result and the
  next member of the list, from left to right.  For instance, if
  \var{action} were the binary addition and \var{value} were 0, then
  folding a list would compute the sum of the list, including for
  empty lists.

\begin{urbiscript}[firstnumber=last]
[].foldl(function (a, b) { a + b }, 0);
[00000000] 0
[1, 2, 3].foldl(function (a, b) { a + b }, 0);
[00000000] 6
[1, 2, 3].foldl(function (a, b) { a - b }, 0);
[00000000] -6
\end{urbiscript}

\item \lstinline|front|\\
  Return the first element of the target. An error if the target is
  empty.

\begin{urbiscript}[firstnumber=last]
assert([0, 1, 2].front == 0);
[].front;
[00000000:error] !!! front: cannot be applied onto empty list
\end{urbiscript}

\item \lstinline|has(\var{that})|\\
  Whether one of the members of the target equals the argument.

\begin{urbiassert}[firstnumber=last]
[0, 1, 2].has(1);
! [0, 1, 2].has(5);
\end{urbiassert}

\item \lstinline|hasSame(\var{that})|\\
  Return whether one of the member of the target is physically equal
  to the argument.

\begin{urbiscript}[firstnumber=last]
var y = 1;
[00000000:hide] 1
[0, y, 2].hasSame(1);
[00000000] false
[0, y, 2].hasSame(y);
[00000000] true
\end{urbiscript}

\item \lstinline|head|\\
  Synonym for \lstinline|front|.

\item \lstinline|insertBack(\var{that})|\\
  Insert the given element at the end of the target.

\begin{urbiscript}[firstnumber=last]
var z = [0, 1];
[00000000] [0, 1]
assert(z.insertBack(2) == [0, 1, 2]);
assert(z == [0, 1, 2]);
\end{urbiscript}

\item \lstinline|insertFront(\var{that})|\\
  Insert the given element at the beginning of the target.

\begin{urbiscript}[firstnumber=last]
var a = [1, 2];
[00000000] [1, 2]
assert(a.insertFront(0) == [0, 1, 2]);
assert(a == [0, 1, 2]);
\end{urbiscript}

\item \lstinline|join(\var{sep} = "", \var{prefix} = "", \var{suffix} = "")|\\
  Bounces to \lstinline|String.join|, see \refObject{String}.

\begin{urbiassert}[firstnumber=last]
["", "ob", ""].join                == "ob";
["", "ob", ""].join("a")           == "aoba";
["", "ob", ""].join("a", "B", "b") == "Baobab";
\end{urbiassert}

\item \lstinline|keys()|\\
  The list of valid indexes.  This allows uniform iteration over a
  \refObject{Dictionary} or a \refObject{List}.

\begin{urbiscript}[firstnumber=last]
{
  var l = ["a", "b", "c"];
  assert(l.keys == [0, 1, 2]);
  assert({
           var res = [];
           for (var k: l.keys)
             res << l[k];
           res
         }
         == l);
};
\end{urbiscript}

\item \lstinline|map(\var{fun})|\\
Apply the given functional value on every member, and return the list
of results.

\begin{urbiassert}[firstnumber=last]
[0, 1, 2, 3].map(function (v) { v % 2 == 0})
        == [true, false, true, false];
\end{urbiassert}

\item \lstinline|range(\var{begin}, \var{end} = nil)|\\
  Return a sub-range of the list, from the first index included to the
  second index excluded.  An error if out of bounds.  Negative indices
  are valid, and number from the end.

  If \var{end} is \lstinline|nil|, calling \lstinline|range(\var{n})
  is equivalent to calling \lstinline|range(0, \var{n})|.

\begin{urbiscript}[firstnumber=last]
do ([0, 1, 2, 3])
{
  assert
  {
    range(0, 0)   == [];
    range(0, 1)   == [0];
    range(1)      == [0];
    range(1, 3)   == [1, 2];

    range(-3, -2) == [1];
    range(-3, -1) == [1, 2];
    range(-3, 0)  == [1, 2, 3];
    range(-3, 1)  == [1, 2, 3, 0];
    range(-4, 4)  == [0, 1, 2, 3, 0, 1, 2, 3];
  };
}|;
[].range(1, 3);
[00428697:error] !!! range: invalid index: 1
\end{urbiscript}

\item \lstinline|remove(\var{val})|\\
  Remove all elements from the target that equals \var{val}.

\begin{urbiscript}[firstnumber=last]
var c = [0, 1, 0, 2, 0, 3];
[00000000] [0, 1, 0, 2, 0, 3]
assert(c.remove(0) == [1, 2, 3]);
assert(c == [1, 2, 3]);
\end{urbiscript}

\item \lstinline|removeBack|\\
  Remove and return the last element of the target. An error if the
  target is empty.

\begin{urbiscript}[firstnumber=last]
var t = [0, 1, 2];
[00000000] [0, 1, 2]
assert(t.removeBack == 2);
assert(t == [0, 1]);
[].removeBack;
[00000000:error] !!! removeBack: cannot be applied onto empty list
\end{urbiscript}

\item \lstinline|removeById(\var{that})|\\
  Remove all elements from the target that physically equals
  \var{that}.

\begin{urbiscript}[firstnumber=last]
var d = 1;
[00000000] 1
var e = [0, 1, d, 1, 2];
[00000000] [0, 1, 1, 1, 2]
assert(e.removeById(d) == [0, 1, 1, 2]);
assert(e == [0, 1, 1, 2]);
\end{urbiscript}

\item \lstinline|removeFront|\\
Remove and return the first element from the target. An error if the
target is empty.

\begin{urbiscript}[firstnumber=last]
var g = [0, 1, 2];
[00000000] [0, 1, 2]
assert(g.removeFront == 0);
assert(g == [1, 2]);
[].removeFront;
[00000000:error] !!! removeFront: cannot be applied onto empty list
\end{urbiscript}

\item \lstinline|reverse|\\
Return the target with the order of elements inverted.

\begin{urbiassert}[firstnumber=last]
[0, 1, 2].reverse == [2, 1, 0];
\end{urbiassert}

\item \lstinline|size|\\
Return the number of elements in the target.

\begin{urbiassert}[firstnumber=last]
[0, 1, 2].size == 3;
[].size == 0;
\end{urbiassert}

\item \lstinline|sort|\\
Return the target, sorted with respect to the \lstinline|<| criteria.

\begin{urbiassert}[firstnumber=last]
[1, 0, 3, 2].sort == [0, 1, 2, 3];
\end{urbiassert}

\item \lstinline|tail|\\
Return the target, minus the first element. An error if the target is
empty.

\begin{urbiscript}[firstnumber=last]
assert([0, 1, 2].tail == [1, 2]);
[].tail;
[00000000:error] !!! tail: cannot be applied onto empty list
\end{urbiscript}

\item \lstinline|'=='(\var{that})|\\
Check whether all elements in the target and \var{that}, are
equal two by two.

\begin{urbiassert}[firstnumber=last]
[0, 1, 2] == [0, 1, 2];
!([0, 1, 2] == [0, 0, 2]);
\end{urbiassert}

\item \lstinline|'[]'(\var{n})|\\
  Return the \var{n}th member of the target (indexing is
  zero-based). If \var{n} is negative, start from the end.  An error
  if out of bounds.

\begin{urbiscript}[firstnumber=last]
assert(["0", "1", "2"][0] == "0");
assert(["0", "1", "2"][2] == "2");
["0", "1", "2"][3];
[00007061:error] !!! []: invalid index: 3

assert(["0", "1", "2"][-1] == "2");
assert(["0", "1", "2"][-3] == "0");
["0", "1", "2"][-4];
[00007061:error] !!! []: invalid index: -4
\end{urbiscript}

\item \lstinline|'[]='(\var{index}, \var{value})|\\
  Assign \var{value} to the element of the target at the given
  \var{index}.

\begin{urbiscript}[firstnumber=last]
var f = [0, 1, 2];
[00000000] [0, 1, 2]
f[1] = 42;
[00000000] 42
assert(f == [0, 42, 2]);
\end{urbiscript}

\item \lstinline|'*'(\var{n})|\\
  Return the target, concatenated \var{n} times to itself.
\begin{urbiassert}[firstnumber=last]
[0, 1] * 3 == [0, 1, 0, 1, 0, 1];
\end{urbiassert}

  Note that since it is the very same list which is repeatedly
  concatenated (the content is not cloned), side-effects on one item
  will reflect on ``all the items''.

\begin{urbiscript}[firstnumber=last]
var l = [[]] * 3;
[00000000] [[], [], []]
l[0] << 1;
[00000000] [1]
l;
[00000000] [[1], [1], [1]]
\end{urbiscript}

\item \lstinline|'+'(\var{other})|\\
Return the concatenation of the target and the \var{other} list.

\begin{urbiassert}[firstnumber=last]
[0, 1] + [2, 3] == [0, 1, 2, 3];
\end{urbiassert}

\item \lstinline|'-'(\var{other})|\\
Return the target without all element that equals any element in the
\var(other) list.

\begin{urbiassert}[firstnumber=last]
[0, 1, 0, 2, 3] - [1, 2] == [0, 0, 3];
\end{urbiassert}

\item \lstinline|'<<'(\var{that})|\\
  A synonym for \lstinline|insertBack|.

\item \lstinline|'<'(\var{other})|\\
  Return whether the target is inferior to the \var{other} list. A
  list is inferior to another if at least one of its element differs
  from the other, and the first differing element is inferior to the
  other.

\begin{urbiassert}[firstnumber=last]
!([0, 1, 2] < [0, 1, 2]);
!([0, 1, 2] < [0, 0, 2]);
[0, 1, 2] < [0, 2, 2];
\end{urbiassert}

  Since List derives from \refObject{Orderable}, the other order-based
  operators are defined.

\begin{urbiassert}[firstnumber=last]
 [0, 1, 2] <= [0, 1, 2];
 [0, 1, 2] >= [0, 1, 2];
 [0, 1, 2] >  [0, 0, 2];
\end{urbiassert}
\end{itemize}

%%% Local Variables:
%%% mode: latex
%%% TeX-master: "../urbi-sdk"
%%% End:

% LocalWords:  lst asList asString foldl hasSame removeBack popback removeFront
% LocalWords:  popfront insertBack pushback insertFront pushfront urbi sameAs
% LocalWords:  removeById setNth


\section{Lobby}
\section{Object}
\label{sect:std-object}

\subsection{Methods}

\subsubsection{createSlot}
\label{sect:std-object-createslot}

\subsubsection{getSlot}
\label{sect:std-object-getslot}

\subsubsection{setSlot}
\label{sect:std-object-setslot}

\subsubsection{updateSlot}
\label{sect:std-object-updateslot}

\section{Orderable}
\label{sect:std-orderable}
\section{Pattern}
\section{Primitive}
\section{String}
\section{Tag}
\section{Tuple}
\section{Void}
\label{sect:std-void}

\chapter{\urbi SDK specifications}

% I just don't know
% what to do
% with myself

\end{document}

% printing, operators, waituntil

% LocalWords:  Urbi monomethods netcat cosinus timestamp subscopes subscope lst
% LocalWords:  getSlot setSlot updateSlot slotNames removeSlot uid arg ok ko ss
% LocalWords:  elt baz quux FIXME getslot lookup protos proto addProto asString
% LocalWords:  removeProto locateSlot locateslot lhs evalArgAt evalArgs ary
% LocalWords:  mytag timeOut myEvent Tuple redefinable CallMessage memoization
% LocalWords:  prepended
